% BHU PYQ -- Professional Exam Template
\documentclass[11pt,a4paper]{article}

\usepackage[a4paper,top=2.5cm,bottom=2.5cm,left=2.8cm,right=2.8cm,headheight=14pt]{geometry}
\usepackage{amsmath,amssymb}
\usepackage[T1]{fontenc}
\usepackage[utf8]{inputenc}
\usepackage{lmodern}
\usepackage{microtype}
\usepackage{fancyhdr}
\usepackage{enumitem}
\usepackage{hyperref}
\usepackage{lastpage}
\usepackage{parskip}

\hypersetup{colorlinks=true,urlcolor=black,linkcolor=black,
  pdftitle={B.Sc. (Hons.) Botany Sem-I BOB-101 -- BHU PYQ}}

\pagestyle{fancy}
\fancyhf{}
\renewcommand{\headrulewidth}{0.4pt}
\renewcommand{\footrulewidth}{0.4pt}
\fancyhead[L]{\small\textit{Banaras Hindu University}}
\fancyhead[R]{\small\textit{Botany Paper: BOB-101}}
\fancyfoot[C]{\small Page \thepage\ of \pageref{LastPage}}

\newlist{parts}{enumerate}{2}
\setlist[parts,1]{label=(\alph*),leftmargin=*,itemsep=4pt,topsep=3pt}
\setlist[parts,2]{label=(\roman*),leftmargin=*,itemsep=2pt,topsep=2pt}
\newcommand{\pts}[1]{\hfill\textit{[#1]}}

\begin{document}

\begin{center}
  {\large\bfseries BANARAS HINDU UNIVERSITY}\\[2pt]
  {\normalsize B.Sc. (Hons.) Semester I Examination, 2023-24}\\[6pt]
  \rule{0.8\linewidth}{0.5pt}\\[6pt]
  {\large\bfseries Botany}\\[4pt]
  {\large\bfseries Paper: BOB-101 --- Cryptogams}\\[4pt]
  {\normalsize Time: Three Hours \hfill Maximum Marks: 70}
\end{center}

\rule{\linewidth}{0.4pt}

\smallskip
\noindent\textit{Instructions: Answer five questions, selecting at least two each from Sections A and B and Section-C which is compulsory. The figures in the right-hand margin indicate marks.}

\bigskip

\begin{center}
  \textbf{Section-A}
\end{center}
\medskip

\medskip\noindent\textbf{Question 1.} Discuss along with suitable diagrams the comparative features of the sporophytes of \textit{Anthoceros} and \textit{Funaria}.\pts{15}

\medskip\noindent\textbf{Question 2.} Write short notes on any two of the following:\pts{$2 \times 7.5 = 15$}
\begin{parts}
  \item Gametophyte of \textit{Marchantia}
  \item False branching in cyanobacteria
  \item Type of zoospores in \textit{Draparnaldiopsis}
\end{parts}

\medskip\noindent\textbf{Question 3.} Discuss the complete life cycle of \textit{Ectocarpus}, along with suitable diagrams.\pts{15}

\medskip\noindent\textbf{Question 4.} Write short notes on any two of the following:\pts{$2 \times 7.5 = 15$}
\begin{parts}
  \item Elaters and pseudoelaters in Bryophytes
  \item Capsule of \textit{Funaria}
  \item Cystocarp
\end{parts}

\bigskip
\begin{center}
  \textbf{Section-B}
\end{center}
\medskip

\medskip\noindent\textbf{Question 5.} Write short notes on any two of the following:\pts{$2 \times 7.5 = 15$}
\begin{parts}
  \item Various types of Steles in Pteridophytes
  \item Outline of Ainsworth's classification
  \item General characters and types of fruiting body in Ascomycotina
\end{parts}

\medskip\noindent\textbf{Question 6.} Write short notes on any two of the following:\pts{$2 \times 7.5 = 15$}
\begin{parts}
  \item Reproduction in \textit{Selaginella}
  \item Reproduction in \textit{Pteris}
  \item Parasexuality in fungi
\end{parts}

\medskip\noindent\textbf{Question 7.} Describe the general characteristic features of Basidiomycotina and describe in brief the life cycle of \textit{Puccinia}.\pts{15}

\medskip\noindent\textbf{Question 8.} Discuss the life cycle of \textit{Equisetum} along with suitable diagrams.\pts{15}

\bigskip
\begin{center}
  \textbf{Section-C}
\end{center}
\medskip

\medskip\noindent\textbf{Question 9.} Answer the following (not more than 50 words each):\pts{$10 \times 1 = 10$}
\begin{parts}
  \item Why do brown algae appear brown in colour?
  \item Who proposed the Stelar theory in Pteridophytes?
  \item Write the name of a Coprophilous fungi.
  \item Dolipore septum is a characteristic feature of which group of fungi?
  \item Pseudoelaters are the characteristic feature of which bryophyte?
  \item Name one algal genus that exhibits stephanokont condition.
  \item Name the sterile tissue formed inside the capsule of \textit{Anthoceros}.
  \item Which structure helps in spore dispersal in \textit{Equisetum}?
  \item Which fungus is commonly known as bread mold?
  \item Gonimoblast filaments are formed in which algal genus?
\end{parts}

\vfill
\rule{\linewidth}{0.4pt}
\begin{center}\small
\href{https://bhu-exam.vercel.app/}{Click Here}\\[4pt]\href{https://me-aryan.vercel.app/}{Click Here}\\[4pt]\href{https://quashan-me.vercel.app/}{Click Here}
\end{center}

\end{document}
